\section{Problem 626 (Round 2)}

\subsection*{1) ROUND-2 OBJECTIVE}

Path C: isolate the exact inverse-threshold obstruction left after Round 1.

Round 1 already reduced the problem to the inverse thresholds
\[
N_k(M):=\min\{n:\exists\,G,\ |V(G)|=n,\ \chi(G)=k,\ \operatorname{girth}(G)>M\}
\]
and
\[
N^{(m)}(t):=\min\{n:\exists\,G,\ |V(G)|=n,\ \operatorname{girth}(G)>m,\ \chi(G)\ge t\},
\]
but only in the one-way form ``if the inverse limit exists, then the original limit exists''.
The natural gap is therefore to prove the exact converse, and more precisely to identify the
\emph{limsup/liminf} of the original quantities in terms of the inverse thresholds.

\subsection*{2) Round-1 FOUNDATION USED}

I use the following Round-1 facts without reproving them.

\begin{itemize}
\item The monotonicity of \(g_k(n)\) and \(h^{(m)}(n)\) in \(n\).
\item The inverse identities
\[
g_k(n)=\max\{M:N_k(M)\le n\},\qquad
h^{(m)}(n)=\max\{t:N^{(m)}(t)\le n\}.
\]
\item The Round-1 one-way transfer statements:
if \(\lim \log N_k(M)/M\) exists then \(\lim g_k(n)/\log n\) exists, and similarly for \(h^{(m)}\).
\end{itemize}

\subsection*{3) NEW INSIGHT / TOOL (ROUND 2)}

The new tool is an exact generalized-inverse theorem at the two scales relevant here:

\begin{enumerate}
\item exponential scale: \(N(M)\approx e^{\lambda M}\) versus \(g(n)\approx (\log n)/\lambda\);
\item power scale: \(N(t)\approx t^{\alpha}\) versus \(h(n)\approx n^{1/\alpha}\).
\end{enumerate}

This upgrades the Round-1 one-way implication to an exact two-sided equivalence, and pins down
the limsup/liminf constants as reciprocal liminf/limsup growth rates of the inverse thresholds.

\subsection*{4) ATTACK PLAN (ROUND 2)}

The exact gaps after Round 1 are:

\begin{itemize}
\item for fixed \(k\ge 4\), Round 1 did not prove that
\[
\lim_{n\to\infty}\frac{g_k(n)}{\log n}
\]
exists if and only if
\[
\lim_{M\to\infty}\frac{\log N_k(M)}{M}
\]
exists;
\item for fixed \(m\ge 3\), Round 1 did not prove the analogous equivalence for
\[
\frac{\log h^{(m)}(n)}{\log n}
\quad\text{and}\quad
\frac{\log N^{(m)}(t)}{\log t}.
\]
\end{itemize}

I close exactly these two gaps by proving reciprocal limsup/liminf formulae.
This overcomes the main obstacle identified in Round 1: we no longer need any extra
regularity assumption beyond monotonicity.

\subsection*{5) WORK (ROUND 2)}

\paragraph{Claim 626.1 (generalized inverse on the exponential scale).}
Let \(N:\mathbb{N}\to\mathbb{N}\) be nondecreasing and satisfy \(N(M)\to\infty\).
Define
\[
g(n):=\max\{M\in\mathbb{N}:N(M)\le n\}.
\]
Put
\[
\lambda_-:=\liminf_{M\to\infty}\frac{\log N(M)}{M},
\qquad
\lambda_+:=\limsup_{M\to\infty}\frac{\log N(M)}{M}.
\]
Assume \(0<\lambda_-\le \lambda_+<\infty\).
Then
\[
\limsup_{n\to\infty}\frac{g(n)}{\log n}=\frac{1}{\lambda_-},
\qquad
\liminf_{n\to\infty}\frac{g(n)}{\log n}=\frac{1}{\lambda_+}.
\]

\emph{Proof.}
Fix \(n\), and let \(M:=g(n)\). By definition of the maximum,
\[
N(M)\le n < N(M+1).
\]
For the \(\limsup\), let \(\varepsilon>0\). By the definition of \(\lambda_-\), for all large \(M\),
\[
\log N(M)\ge (\lambda_--\varepsilon)M.
\]
Therefore, for all large \(n\),
\[
\log n\ge \log N(M)\ge (\lambda_--\varepsilon)M,
\]
so
\[
\frac{g(n)}{\log n}=\frac{M}{\log n}\le \frac{1}{\lambda_--\varepsilon}.
\]
Hence
\[
\limsup_{n\to\infty}\frac{g(n)}{\log n}\le \frac{1}{\lambda_-}.
\]

For the reverse inequality, choose a sequence \(M_j\to\infty\) such that
\[
\frac{\log N(M_j)}{M_j}\to \lambda_-.
\]
Set \(n_j:=N(M_j)\). Then \(g(n_j)\ge M_j\), so
\[
\frac{g(n_j)}{\log n_j}\ge \frac{M_j}{\log N(M_j)}\to \frac{1}{\lambda_-}.
\]
Thus
\[
\limsup_{n\to\infty}\frac{g(n)}{\log n}\ge \frac{1}{\lambda_-}.
\]
This proves the first identity.

Now turn to the \(\liminf\). Let \(\varepsilon>0\). By the definition of \(\lambda_+\), for all large \(M\),
\[
\log N(M)\le (\lambda_++\varepsilon)M.
\]
Using \(n<N(M+1)\), we get
\[
\log n<\log N(M+1)\le (\lambda_++\varepsilon)(M+1).
\]
Hence
\[
\frac{g(n)}{\log n}=\frac{M}{\log n}
>
\frac{M}{(\lambda_++\varepsilon)(M+1)}.
\]
Since \(M=g(n)\to\infty\) as \(n\to\infty\), it follows that
\[
\liminf_{n\to\infty}\frac{g(n)}{\log n}\ge \frac{1}{\lambda_++\varepsilon}.
\]
Letting \(\varepsilon\downarrow 0\) gives
\[
\liminf_{n\to\infty}\frac{g(n)}{\log n}\ge \frac{1}{\lambda_+}.
\]

For the reverse inequality, choose \(M_j\to\infty\) with
\[
\frac{\log N(M_j)}{M_j}\to \lambda_+.
\]
Set \(n_j:=N(M_j)-1\). Since \(N(M_j)\to\infty\), we have \(n_j\to\infty\), and because
\(n_j<N(M_j)\), we get \(g(n_j)\le M_j-1\). Therefore
\[
\frac{g(n_j)}{\log n_j}\le
\frac{M_j-1}{\log(N(M_j)-1)}.
\]
Now \(\log(N(M_j)-1)\sim \log N(M_j)\), so
\[
\frac{M_j-1}{\log(N(M_j)-1)}\to \frac{1}{\lambda_+}.
\]
Hence
\[
\liminf_{n\to\infty}\frac{g(n)}{\log n}\le \frac{1}{\lambda_+}.
\]
This proves the second identity. \(\square\)

\paragraph{Claim 626.2 (application to \(g_k(n)\)).}
For fixed \(k\ge 4\), define
\[
\lambda_k^-:=\liminf_{M\to\infty}\frac{\log N_k(M)}{M},
\qquad
\lambda_k^+:=\limsup_{M\to\infty}\frac{\log N_k(M)}{M}.
\]
Then
\[
\limsup_{n\to\infty}\frac{g_k(n)}{\log n}=\frac{1}{\lambda_k^-},
\qquad
\liminf_{n\to\infty}\frac{g_k(n)}{\log n}=\frac{1}{\lambda_k^+}.
\]

\emph{Proof.}
Round 1 already established the inverse identity
\[
g_k(n)=\max\{M:N_k(M)\le n\}.
\]
Also \(N_k(M)\to\infty\): indeed, any graph of girth \(>M\) containing a cycle must have at least \(M+1\) vertices, and \(\chi(G)=k\ge 4\) certainly forces \(G\) to contain a cycle.
Thus Claim 626.1 applies directly. \(\square\)

\paragraph{Corollary 626.3 (exact equivalence for the first limit).}
For fixed \(k\ge 4\),
\[
\lim_{n\to\infty}\frac{g_k(n)}{\log n}
\quad\text{exists}
\quad\Longleftrightarrow\quad
\lim_{M\to\infty}\frac{\log N_k(M)}{M}
\quad\text{exists},
\]
and in that case the two limits are reciprocal.

\emph{Proof.}
The forward and backward implications follow immediately from Claim 626.2, because a real sequence has a limit if and only if its limsup and liminf are equal. \(\square\)

\paragraph{Claim 626.4 (generalized inverse on the power scale).}
Let \(N:\mathbb{N}\to\mathbb{N}\) be nondecreasing and \(N(t)\to\infty\).
Define
\[
h(n):=\max\{t\in\mathbb{N}:N(t)\le n\}.
\]
Put
\[
\alpha_-:=\liminf_{t\to\infty}\frac{\log N(t)}{\log t},
\qquad
\alpha_+:=\limsup_{t\to\infty}\frac{\log N(t)}{\log t}.
\]
Assume \(0<\alpha_-\le \alpha_+<\infty\).
Then
\[
\limsup_{n\to\infty}\frac{\log h(n)}{\log n}=\frac{1}{\alpha_-},
\qquad
\liminf_{n\to\infty}\frac{\log h(n)}{\log n}=\frac{1}{\alpha_+}.
\]

\emph{Proof.}
Fix \(n\), and let \(t:=h(n)\). By definition,
\[
N(t)\le n < N(t+1).
\]
For the \(\limsup\), let \(\varepsilon>0\). By the definition of \(\alpha_-\), for all large \(t\),
\[
\log N(t)\ge (\alpha_--\varepsilon)\log t.
\]
Since \(n\ge N(t)\), we get
\[
\log n\ge \log N(t)\ge (\alpha_--\varepsilon)\log t,
\]
so
\[
\frac{\log h(n)}{\log n}=\frac{\log t}{\log n}\le \frac{1}{\alpha_--\varepsilon}.
\]
Thus
\[
\limsup_{n\to\infty}\frac{\log h(n)}{\log n}\le \frac{1}{\alpha_-}.
\]

For the reverse inequality, choose \(t_j\to\infty\) such that
\[
\frac{\log N(t_j)}{\log t_j}\to \alpha_-.
\]
Set \(n_j:=N(t_j)\). Then \(h(n_j)\ge t_j\), hence
\[
\frac{\log h(n_j)}{\log n_j}\ge \frac{\log t_j}{\log N(t_j)}\to \frac{1}{\alpha_-}.
\]
Therefore
\[
\limsup_{n\to\infty}\frac{\log h(n)}{\log n}\ge \frac{1}{\alpha_-}.
\]
This proves the first identity.

Now let \(\varepsilon>0\). By the definition of \(\alpha_+\), for all large \(t\),
\[
\log N(t)\le (\alpha_++\varepsilon)\log t.
\]
Using \(n<N(t+1)\), we obtain
\[
\log n<\log N(t+1)\le (\alpha_++\varepsilon)\log(t+1),
\]
so
\[
\frac{\log h(n)}{\log n}=\frac{\log t}{\log n}>
\frac{\log t}{(\alpha_++\varepsilon)\log(t+1)}.
\]
As \(t=h(n)\to\infty\), the right-hand side tends to \(1/(\alpha_++\varepsilon)\). Hence
\[
\liminf_{n\to\infty}\frac{\log h(n)}{\log n}\ge \frac{1}{\alpha_++\varepsilon}.
\]
Letting \(\varepsilon\downarrow 0\) yields
\[
\liminf_{n\to\infty}\frac{\log h(n)}{\log n}\ge \frac{1}{\alpha_+}.
\]

For the reverse inequality, choose \(t_j\to\infty\) with
\[
\frac{\log N(t_j)}{\log t_j}\to \alpha_+,
\]
and set \(n_j:=N(t_j)-1\). Then \(h(n_j)\le t_j-1\), so
\[
\frac{\log h(n_j)}{\log n_j}
\le
\frac{\log(t_j-1)}{\log(N(t_j)-1)}.
\]
Since \(\log(t_j-1)\sim \log t_j\) and \(\log(N(t_j)-1)\sim \log N(t_j)\), the right-hand side tends to \(1/\alpha_+\). Therefore
\[
\liminf_{n\to\infty}\frac{\log h(n)}{\log n}\le \frac{1}{\alpha_+}.
\]
This proves the second identity. \(\square\)

\paragraph{Claim 626.5 (application to \(h^{(m)}(n)\)).}
For fixed \(m\ge 3\), define
\[
\alpha_m^-:=\liminf_{t\to\infty}\frac{\log N^{(m)}(t)}{\log t},
\qquad
\alpha_m^+:=\limsup_{t\to\infty}\frac{\log N^{(m)}(t)}{\log t}.
\]
Then
\[
\limsup_{n\to\infty}\frac{\log h^{(m)}(n)}{\log n}=\frac{1}{\alpha_m^-},
\qquad
\liminf_{n\to\infty}\frac{\log h^{(m)}(n)}{\log n}=\frac{1}{\alpha_m^+}.
\]

\emph{Proof.}
Round 1 already established
\[
h^{(m)}(n)=\max\{t:N^{(m)}(t)\le n\}.
\]
Also \(N^{(m)}(t)\to\infty\), since any graph with chromatic number at least \(t\) has at least \(t\) vertices. So Claim 626.4 applies directly. \(\square\)

\paragraph{Corollary 626.6 (exact equivalence for the second limit).}
For fixed \(m\ge 3\),
\[
\lim_{n\to\infty}\frac{\log h^{(m)}(n)}{\log n}
\quad\text{exists}
\quad\Longleftrightarrow\quad
\lim_{t\to\infty}\frac{\log N^{(m)}(t)}{\log t}
\quad\text{exists},
\]
and in that case the two limits are reciprocal.

\emph{Proof.}
Again this is immediate from Claim 626.5. \(\square\)

\paragraph{Corollary 626.7 (translation of the currently listed bounds to the inverse scale).}
Using the current bounds recorded on the problem page,
\[
\frac{1}{4\log k}\log n\le g_k(n)\le \frac{2}{\log(k-2)}\log n+1,
\]
one gets
\[
\frac{\log(k-2)}{2}
\le
\liminf_{M\to\infty}\frac{\log N_k(M)}{M}
\le
\limsup_{M\to\infty}\frac{\log N_k(M)}{M}
\le
4\log k.
\]

\emph{Proof.}
Let \(a:=1/(4\log k)\) and \(b:=2/\log(k-2)\).

From \(g_k(n)\ge a\log n\), choose \(n=\lceil e^{M/a}\rceil\). Then \(g_k(n)\ge M\), hence
\(N_k(M)\le n\), so
\[
\limsup_{M\to\infty}\frac{\log N_k(M)}{M}\le \frac{1}{a}=4\log k.
\]

From \(g_k(n)\le b\log n+1\), if \(n<e^{(M-1)/b}\) then \(g_k(n)<M\), so
\(N_k(M)\ge e^{(M-1)/b}\). Therefore
\[
\liminf_{M\to\infty}\frac{\log N_k(M)}{M}\ge \frac{1}{b}=\frac{\log(k-2)}{2}.
\]
This proves the claim. \(\square\)

\subsection*{6) ADVERSARIAL VERIFICATION}

The new work was checked against the following possible failure modes.

\begin{itemize}
\item \emph{Non-strict growth of \(N_k(M)\) and \(N^{(m)}(t)\).}  
The proofs only use monotonicity, not strict monotonicity. The identities
\[
N(g(n))\le n < N(g(n)+1)
\]
and
\[
N(h(n))\le n < N(h(n)+1)
\]
remain valid for generalized inverses.

\item \emph{Use of \(N(M)-1\).}  
This requires \(N(M)\ge 2\), which holds for all sufficiently large \(M\) because \(N(M)\to\infty\).

\item \emph{Could \(g_k(n)\) or \(h^{(m)}(n)\) stay bounded?}  
No. If \(g_k(n)\) were bounded along an unbounded sequence of \(n\), then \(N_k(M)\) would be bounded for some arbitrarily large \(M\), impossible. The same argument works for \(h^{(m)}\).

\item \emph{Reciprocal-limsup/lower-limsup interchange.}  
Each direction was proved separately; no hidden continuity assumption was used.
\end{itemize}

\subsection*{7) FINAL}

\textbf{UNRESOLVED (BUT STRICTLY ADVANCED).}

The strict advance over Round 1 is the exact reciprocal limsup/liminf theory:
the two Erd\H{o}s limits in Problem 626 exist \emph{if and only if} the corresponding inverse-threshold
limits exist, and the limsup/liminf constants are exactly reciprocals of the inverse liminf/limsup.
Thus the problem is now reduced with no gap to determining the asymptotic growth of
\(N_k(M)\) and \(N^{(m)}(t)\).

\subsection*{8) COMPLETION ESTIMATE}

\textbf{COMPLETION: 40\%}

\subsection*{9) REFERENCES}

\begin{itemize}
\item[\textbf{[Bloom626]}] T. F. Bloom, \emph{Erd\H{o}s Problem \#626}, accessed 2026-03-08.
\end{itemize}

